\documentclass[a4paper,11pt,lualatex,ja=standard]{bxjsarticle}
\geometry{top=20mm,bottom=20mm,left=25mm,right=25mm}
\usepackage{graphicx}
\usepackage{booktabs}
\usepackage{array}
\usepackage{xcolor}

% 行間を詰める
\setlength{\parskip}{0pt}
\setlength{\baselineskip}{18pt}

% 3文書共通の体裁設定（和文フォント・見出しの番号様式）
% 研究計画書（kyodai-protocol.cls）と同じファイルを読み込む
% ===================================================================
% 3文書（研究計画書・説明文書・同意書）で共有する体裁設定
%
% 京都大学 医の倫理委員会「書類を作成する際のチェックリスト」の
%   ・使用されるフォントは「MSゴシック」「MS明朝」としてください
%   ・インデント、フォントサイズ、項目番号の表記の体裁を整えてください
% に対応するため、和文フォントと見出しの番号様式をこのファイルに集約する。
% kyodai-protocol.cls、explanation_template.tex、consent_template.tex の
% いずれからも読み込まれるため、3文書の体裁が自動的に揃う。
% ===================================================================

% -------------------------------------------------------------------
% 和文フォント: IPA明朝 / IPAゴシック
%
% チェックリストは「MS明朝」「MSゴシック」を指定するが、当該フォントは
% Microsoft専有でありビルド環境（texlive/texlive イメージ）に導入できない。
% 規定の趣旨である「Windowsとの互換性による文字化けの防止」はPDFへの
% フォント埋め込みによって満たしたうえで、書体としては指定に最も近い
% 無償の標準和文フォント（IPA明朝／IPAゴシック）を用いる。
% -------------------------------------------------------------------
\RequirePackage[ipa]{luatexja-preset}

% -------------------------------------------------------------------
% 見出しの番号様式
%
% 医の倫理委員会が示す「研究計画書に記載すべき事項」および実際に承認された
% 計画書で用いられている表記に揃える。
%   section       → 1.  2.  3. …
%   subsection    → 1)  2)  3) …
%   subsubsection → ①  ②  ③ …
%   paragraph     → 番号なし（太字見出し）
% -------------------------------------------------------------------
\RequirePackage{titlesec}

\newcommand{\jpcirclednum}[1]{%
  \ifcase#1\or ①\or ②\or ③\or ④\or ⑤\or ⑥\or ⑦\or ⑧\or ⑨\or ⑩%
  \or ⑪\or ⑫\or ⑬\or ⑭\or ⑮\or ⑯\or ⑰\or ⑱\or ⑲\or ⑳\else (\number#1)\fi}

\renewcommand{\thesection}{\arabic{section}.}
\renewcommand{\thesubsection}{\arabic{subsection})}
\renewcommand{\thesubsubsection}{\jpcirclednum{\value{subsubsection}}}

% 下位の見出しほど左を字下げし、階層を視覚的に示す
\titleformat{\section}[hang]{\normalfont\Large\bfseries}{\thesection}{0.7em}{}
\titlespacing*{\section}{0pt}{3.5ex plus 1ex minus .2ex}{2.3ex plus .2ex}
\titleformat{\subsection}[hang]{\normalfont\large\bfseries}{\thesubsection}{0.6em}{}
\titlespacing*{\subsection}{1em}{3.0ex plus 1ex minus .2ex}{1.8ex plus .2ex}
\titleformat{\subsubsection}[hang]{\normalfont\normalsize\bfseries}{\thesubsubsection}{0.4em}{}
\titlespacing*{\subsubsection}{2em}{2.6ex plus 1ex minus .2ex}{1.4ex plus .2ex}
\titleformat{\paragraph}[hang]{\normalfont\normalsize\bfseries}{}{0em}{}
\titlespacing*{\paragraph}{3em}{2.2ex plus 0.8ex minus .2ex}{1.0ex plus .2ex}

% 番号を付す深さはsubsubsection（①）まで。paragraphは番号なしの太字見出し。
\setcounter{secnumdepth}{3}


% 研究の基本情報を読み込み（protocol_template.texと共通）
\input{research_info.tex}

\begin{document}

\begin{center}
{\LARGE \textbf{研究参加の同意書}}

\vspace{0.5em}

{\small \ProtocolVersion\hspace{1em}\ProtocolDate}
\end{center}

\vspace{1em}

\noindent
【研究を行う機関の長：（\InstitutionHead）等】殿

\vspace{0.5em}

\noindent
説明を受け、以下の項目を理解し、【\ResearchTitle】に参加することに同意します。

\vspace{0.5em}

\begin{enumerate}
\setlength{\itemsep}{0.1em}
\setlength{\parsep}{0pt}
\item 研究実施について
\item 研究機関
\item 研究の目的および意義
\item 研究方法と期間
\item 研究対象者として選定された理由
\item 研究対象者に生じる負担と、予想されるリスクおよび利益
\item いつでも同意の撤回ができます
\item 研究に同意しない、または同意撤回において不利益はないこと
\item 研究に関する情報公開について
\item 研究計画書等の閲覧について
\item 個人情報の取扱いについて
\item 試料・情報の保管および廃棄の方法
\item 研究資金および利益相反について
\item 研究より得られた結果の取り扱い
\item 相談の窓口
\item 経済的負担／謝礼について
% 該当する場合のみ以下の項目を含めてください
% \item （通常の診療を超える医療行為を伴う研究の場合）他の治療法について
% \item （通常の診療を超える医療行為を伴う研究の場合）研究参加後の医療の提供について
% \item （侵襲を伴う研究の場合）健康被害の補償について
\item 試料・情報の将来の研究における使用および他機関への提供
% \item （侵襲を伴う研究であって介入を行うものの場合）モニタリング・監査・倫理審査委員会による個人情報の閲覧
\end{enumerate}

% \vspace{0.5em}
% {\small \textcolor{red}{\textbf{※省略した項目がある場合、項目の通し番号に注意してください。}}}

\vspace{1.5em}

\begin{tabular}{ll}
日付 & （西暦）\hspace{3em}年\hspace{2em}月\hspace{2em}日 \\[1.5em]
本人氏名（自筆） & \hspace{10cm} \\[2em]
\end{tabular}

\vspace{1em}

\noindent
（代諾者が必要であれば）

\vspace{0.5em}

\begin{tabular}{ll}
代諾者 氏名（自筆） & \hspace{6cm} 本人との関係（\hspace{3cm}） \\[2em]
\end{tabular}

\vspace{1em}

\begin{tabular}{ll}
説明者 氏名（自筆） & \hspace{6cm} （所属）\hspace{3cm} \\[1.5em]
\end{tabular}

\vspace{1em}

\noindent
{\small \textbf{※同意書のコピーを一部必ずもらってください}}

\end{document}
